% ---------------------------------------------------------------------
% Related Work (~2 pages). Structured to demonstrate coverage of
% 2024--2026 literature (surveys, foundation models, Mamba, ALLMs).
% ---------------------------------------------------------------------

This section situates AuralGuard within five research streams: (i)~the evolution of
spoofing countermeasures from hand-crafted features to foundation-model detectors,
(ii)~self-supervised representations and layer-aggregation strategies,
(iii)~one-class and open-set objectives, (iv)~generalization and robustness studies,
and (v)~the emerging role of large audio-language models. Recent
surveys~\cite{yi2023surveyadd,pham2025surveydeepfake,li2025survey} provide broader
coverage; Table~\ref{tab:positioning} positions our work against the most closely
related recent systems.

\subsection{Spoofing Countermeasures: Four Generations}
\label{ssec:related-cm}

\textbf{Generation 1 --- hand-crafted features.} Early ASVspoof systems paired
cepstral features with Gaussian mixture models or shallow classifiers. Constant-Q
cepstral coefficients~\cite{todisco2017constantq} and linear-frequency cepstral
coefficients~\cite{sahidullah2015comparison} remain reference front-ends, and
LFCC--GMM is still reported as an official challenge baseline. These systems saturate
on modern synthesis and transfer poorly across corpora~\cite{wang2020asvspoof2019}.

\textbf{Generation 2 --- end-to-end deep networks.} Light CNNs over cepstral
inputs~\cite{lavrentyeva2019stc} were followed by raw-waveform models:
RawNet2~\cite{tak2021rawnet2} learns SincConv filterbanks directly from audio, and
AASIST~\cite{jung2022aasist} introduced a heterogeneous spectro-temporal graph
attention back-end that models spectral and temporal artifact patterns jointly,
becoming the standard back-end for subsequent work.

\textbf{Generation 3 --- SSL foundation-model detectors.} Tak et
al.~\cite{tak2022automatic} showed that a wav2vec~2.0 front-end coupled with AASIST
and RawBoost augmentation~\cite{tak2022rawboost} reduces in-domain EER several-fold,
establishing the SSL+AASIST recipe that dominates current leaderboards. Variants
substitute WavLM~\cite{chen2022wavlm} or XLS-R~\cite{babu2022xlsr} front-ends, add
KAN-enhanced heads~\cite{borodin2024aasist3}, or attach sensitive-layer-selection
classifiers~\cite{zhang2024slsxlsr}.

\textbf{Generation 4 --- post-transformer architectures (2024--2025).} The most recent
wave replaces or augments the transformer/graph pipeline: RawBMamba~\cite{chen2024rawbmamba}
and XLSR-Mamba~\cite{xiao2025xlsrmamba} exploit bidirectional state-space models for
long-range artifact modeling, with XLSR-Mamba reporting state-of-the-art results on
ASVspoof~2021 and In-the-Wild; Nes2Net~\cite{liu2025nes2net} shows that a lightweight
nested architecture can replace heavy back-ends on top of foundation-model features;
TCM~\cite{truong2024tcm} models temporal--channel dependencies inside multi-head
attention; and SLIM~\cite{zhu2024slim} detects deepfakes through an explicit
style--linguistics mismatch prior, improving out-of-domain transfer. SafeEar~\cite{li2024safeear}
addresses the orthogonal problem of privacy-preserving detection from decoupled
acoustic tokens. AuralGuard belongs to this generation but differs in that it
\emph{combines} a foundation-model view with an explicit low-level artifact view,
rather than relying on either alone.

\subsection{SSL Representations and Layer Aggregation}
\label{ssec:related-ssl}

Speech foundation models---wav2vec~2.0~\cite{baevski2020wav2vec},
HuBERT~\cite{hsu2021hubert}, WavLM~\cite{chen2022wavlm}, XLS-R~\cite{babu2022xlsr},
and Whisper~\cite{radford2023whisper}---encode a hierarchy in which lower transformer
layers capture acoustic detail and upper layers capture phonetic and semantic content.
Where spoofing artifacts surface within this hierarchy varies by generator, which
makes the aggregation strategy a first-order design decision. Prior detectors use the
final layer~\cite{tak2022automatic}, a fixed weighted sum, or concatenation of the
last~$k$ layers; Pan et al.~\cite{pan2024attentive} recently showed that attentive
merging of hidden embeddings measurably improves anti-spoofing, and Pascu et
al.~\cite{pascu2024calibrated} found that even \emph{frozen} SSL representations with
simple aggregation generalize surprisingly well when properly calibrated. AuralGuard
adopts learnable softmax layer attention with an entropy regularizer that prevents
collapse onto a single layer, and we analyze the learned layer profile
(Section~\ref{ssec:layeranalysis}) to show it concentrates on lower-middle layers
where vocoder artifacts reside.

\subsection{One-Class and Open-Set Objectives}
\label{ssec:related-oneclass}

Binary cross-entropy treats spoof detection as closed-set discrimination between
bona-fide speech and \emph{known} attacks; unseen attacks at test time may fall
anywhere relative to the learned boundary. One-class formulations instead model the
bona-fide distribution compactly and score outlierness.
OC-Softmax~\cite{zhang2021oneclass} imposes asymmetric angular margins around a
learnable bona-fide center and remains the strongest simple open-set objective; recent
extensions include adaptive centroid shift~\cite{kim2024oneclassknowledge} and
continual-learning variants that adapt the one-class region as new attacks
emerge~\cite{zhang2024remember}. Supervised contrastive learning~\cite{khosla2020supcon}
complements these margins by structuring the embedding neighborhood. AuralGuard
combines OC-Softmax with a SupCon auxiliary and shows in ablation that the pairing
outperforms either objective alone under domain shift.

\subsection{Generalization, Robustness, and Benchmarks}
\label{ssec:related-gen}

M{\"u}ller et al.~\cite{muller2022inthewild} first quantified the collapse of
laboratory detectors on found-in-the-wild audio, and their follow-up
study~\cite{muller2024harder} attributes cross-domain failure less to ``harder''
attacks than to systematically \emph{different} artifact distributions---an argument
for multi-view features and distribution-tightening objectives such as ours. Latent
space refinement and augmentation~\cite{wang2025latentrefine} tackle the same gap at
the representation level. On the data side, the field has moved quickly:
ASVspoof~5~\cite{wang2024asvspoof5,wang2025asvspoof5journal} introduces crowdsourced
speech, surrogate-optimized adversarial attacks, and a deepfake track;
MLAAD~\cite{muller2024mlaad} spans 38 languages and 82 TTS systems;
Codecfake~\cite{xie2025codecfake} targets neural-codec-based synthesis;
SpoofCeleb~\cite{jung2025spoofceleb} sources training data from real-world VoxCeleb
conditions; and domain-specific corpora address singing voice~\cite{zang2024singfake}
and LLM-scripted disinformation~\cite{luong2024llamapartialspoof}. Channel robustness
was isolated as a primary axis by ASVspoof~2021~\cite{liu2023asvspoof2021};
RawBoost~\cite{tak2022rawboost} remains the standard signal-level augmentation, which
our curriculum extends with a realistic transcoding chain, measured room impulse
responses~\cite{ko2017rir}, and MUSAN noise~\cite{snyder2015musan}.

\subsection{Large Audio-Language Models and Deepfake Detection}
\label{ssec:related-allm}

Instruction-following audio-language models such as Qwen2-Audio~\cite{chu2024qwen2audio}
and SALMONN~\cite{tang2024salmonn} unify perception and reasoning over audio, and it
is natural to ask whether such general-purpose models make specialized countermeasures
obsolete. Current evidence indicates the opposite: ALLM4ADD~\cite{gong2025allm4add}
reports that leading ALLMs prompted zero-shot perform close to chance on deepfake
detection, and only approach competitive accuracy after task-specific supervised
fine-tuning that reframes detection as audio question answering. However, published
comparisons evaluate ALLMs on single corpora under clean conditions. Our protocol is,
to our knowledge, the first to place zero-shot ALLM prompting alongside four
generations of specialized countermeasures under identical cross-domain,
channel-distorted, and calibration-aware evaluation, quantifying precisely where the
foundation-model generality does and does not help.

\subsection{Positioning}
\label{ssec:related-positioning}

\begin{table}[t]
\centering
\caption{Positioning of AuralGuard against closely related recent detectors.
  SSL~= foundation-model front-end; Art.~= explicit low-level artifact features;
  OC~= one-class objective; Chan.~= channel/codec augmentation;
  Cal.~= calibration reported; X-ling.~= cross-lingual evaluation.}
\label{tab:positioning}
\setlength{\tabcolsep}{3.2pt}
\footnotesize
\begin{tabular}{lcccccc}
\toprule
System & SSL & Art. & OC & Chan. & Cal. & X-ling. \\
\midrule
W2V2+AASIST~\cite{tak2022automatic}      & \checkmark & --- & --- & \checkmark & --- & --- \\
OC-Softmax~\cite{zhang2021oneclass}      & --- & \checkmark & \checkmark & --- & --- & --- \\
SLIM~\cite{zhu2024slim}                  & \checkmark & --- & --- & --- & --- & --- \\
XLSR-Mamba~\cite{xiao2025xlsrmamba}      & \checkmark & --- & --- & \checkmark & --- & --- \\
Nes2Net~\cite{liu2025nes2net}            & \checkmark & --- & --- & \checkmark & --- & \checkmark \\
Pascu et al.~\cite{pascu2024calibrated}  & \checkmark & --- & --- & --- & \checkmark & --- \\
ALLM4ADD~\cite{gong2025allm4add}         & \checkmark & --- & --- & --- & --- & --- \\
\midrule
\textbf{AuralGuard (ours)} & \checkmark & \checkmark & \checkmark & \checkmark & \checkmark & \checkmark \\
\bottomrule
\end{tabular}
\end{table}

In summary, AuralGuard is distinguished from prior work by the \emph{combination} of
(i)~multi-view learning that pairs a foundation-model stream with an explicit
phase/group-delay artifact stream, (ii)~gated cross-attention fusion that weights the
views per input, (iii)~a one-class contrastive objective with layer-attention entropy
regularization, and (iv)~an evaluation protocol that reports generalization,
robustness, cross-lingual transfer, calibration, and explainability together
(Table~\ref{tab:positioning}). No prior system covers all six columns of
Table~\ref{tab:positioning} simultaneously.
